% ============================================
% 选择题 Q11-20 讲稿
% ============================================
\section{选择题（二）}

% -------------------------------------------------------
\subsection{第11题：函数形参}
\textbf{概念概要：} 函数形参的语法规则——形参的类型、void 函数、形参与全局变量重名。

\textbf{核心讲解要点：}
\begin{itemize}
    \item A 正确：形参可以是任何类型——整型、指针、结构体、数组（退化为指针）等都合法。让学生回忆见过的各种函数声明。
    \item B 错误：\texttt{void} 函数完全可以有形参，例如 \texttt{void printf\_stars(int n);}。\texttt{void} 只是”不返回值”，与形参无关。
    \item C 错误：形参名会”屏蔽”同名的全局变量。这是 C 语言的\textbf{名字遮蔽}机制——局部作用域优先。可以举例：在函数体内你无法访问被屏蔽的全局变量。
    \item D 错误：无形参函数示例——\texttt{int getch(void);}、\texttt{int rand(void);}。
\end{itemize}

\textbf{常见错误：}
\begin{itemize}
    \item 学生常以为 \texttt{void func()} 就是”不能有参数”，实际上 \texttt{void} 只修饰返回类型。
    \item 混淆”形参名遮蔽全局变量”和”不能同名”——语法上允许同名，结果是屏蔽而非冲突。
\end{itemize}

\textbf{讲师提示：} \textit{此题为送分题，快速过。重点强调”形参能屏蔽全局变量”这个知识点，后面的阅读程序可能会遇到。}

{\color{highlight}答案：A}

% -------------------------------------------------------
\subsection{第12题：宏定义展开陷阱}

\textbf{概念概要：} 宏是\underline{纯文本替换}，不进行任何计算。这是宏与函数最本质的区别。

\smallskip
\textbf{讲师提示：} \textit{这是整个选择题中最重要的坑！务必让学生亲手在纸上展开一遍。}

\textbf{展开过程 — 带学生逐步完成：}

\noindent 宏定义：\texttt{\#define F(r)\quad r+r}

\noindent 表达式：\texttt{F(3) * F(3)}

\noindent \textbf{第1步}：将实参 3 替换形参 r，得到 \texttt{3+3 * 3+3}

\noindent \textbf{第2步}：\underline{此时才进行运算符优先级分析}——乘法优先于加法

\noindent \textbf{第3步}：$3 + (3 \times 3) + 3 = 3 + 9 + 3 = \mathbf{15}$

\bigskip
\textbf{为什么不是 36？} 学生直觉往往会”先算 \texttt{F(3)} 得 6，再 $6 \times 6 = 36$”。但宏不先计算——它只做文本粘贴。要想得到 36，宏定义必须加括号：

\begin{lstlisting}[language=C]
#define F(r) ((r)+(r))   // 安全的宏定义
\end{lstlisting}

\noindent 此时展开为 \texttt{((3)+(3)) * ((3)+(3)) = 6 * 6 = 36}。

\textbf{常见错误：}
\begin{itemize}
    \item 把宏当函数用，以为会”先求值再替换”——这是最根本的错误。
    \item 忘记宏的副作用：\texttt{F(i++)} 展开为 \texttt{i++ + i++}，i 自增两次，行为未定义！
    \item 不加括号导致的优先级问题不仅限于算术，逻辑运算符也一样（如 \texttt{\#define AND(a,b) a\&\&b} 写成 \texttt{AND(x, y)==1} 的歧义）。
\end{itemize}

{\color{highlight}答案：C（15）}

% -------------------------------------------------------
\subsection{第13题：函数的定义与调用嵌套}

\textbf{概念概要：} C 语言函数的定义与调用的嵌套规则。

\textbf{核心讲解要点：}
\begin{itemize}
    \item \textbf{定义不可嵌套}：C 语言中所有函数定义都是”平等”的，不允许在一个函数体内定义另一个函数。这与 Pascal、Python 等语言不同。C 没有”内部函数”的概念。
    \item \textbf{调用可以嵌套}：函数 A 调用函数 B，B 调用函数 C——这是标准的嵌套调用。递归更是一种特殊的嵌套：函数调用自身。
    \item 标准 C（C89/C99/C11）不支持嵌套函数定义。GCC 有扩展支持但\underline{不可移植}，考试中视为非法。
\end{itemize}

\textbf{常见错误：}
\begin{itemize}
    \item 有 Python 基础的学生可能习惯在函数内定义辅助函数，会错误认为 C 也支持。
    \item 混淆”嵌套调用”与”递归调用”——递归是调用的特例，”嵌套调用”更宽泛。
\end{itemize}

{\color{highlight}答案：B（定义不可嵌套，调用可嵌套）}

% -------------------------------------------------------
\subsection{第14题：指针下标运算}

\textbf{概念概要：} 指针的下标运算本质是指针算术——\texttt{p[i] $\equiv$ *(p + i)}，\texttt{p[-3] $\equiv$ *(p - 3)}。

\textbf{核心讲解要点：}
\begin{itemize}
    \item 给定 \texttt{int *p = \&a[5];}，即 p 指向数组元素 a[5]（值为 6）。
    \item \texttt{p[-3]} 等价于 \texttt{*(p - 3)}：从 a[5] 向低地址方向移动 3 个 int 元素，到达 a[2]。
    \item a[2] 的值为 \textbf{3}。
    \item \textbf{关键理解：} 数组下标\underline{本质就是指针偏移}。a[i] 不过是 *(a+i) 的语法糖。因此下标可以为负，只要能合法访问即可。
\end{itemize}

\noindent 图示（帮助学生可视化）：

\begin{center}
\texttt{a: [1] [2] [3] [4] [5] [6] [7] [8] [?] [?]}\\
\texttt{\qquad\qquad\ \^{}p[-3] \qquad\qquad\ \^{}p}
\end{center}

\textbf{常见错误：}
\begin{itemize}
    \item 认为”下标不能是负数”——这是对数组下标的误解。下标只是偏移量。
    \item 把 p[-3] 理解为”从 p 往回数 3 个”是对的，但计算具体是哪个元素时容易数错（\texttt{p} 指向 a[5]，p[-1] → a[4]，p[-2] → a[3]，p[-3] → a[2]）。
\end{itemize}

\textbf{讲师提示：} \textit{花 30 秒画一个数组内存布局图，让学生直观感受 p-3 移动的位置。强调 “下标就是偏移量” 这个本质。}

{\color{highlight}答案：C（3）}

% -------------------------------------------------------
\subsection{第15题：联合体（union）大小}

\textbf{概念概要：} union 的所有成员\textbf{共享同一块内存空间}，大小等于最大成员的大小。

\textbf{核心讲解要点：}
\begin{itemize}
    \item \texttt{unionexp\{int i; float j; double k;\} x;} 中：int 占 4 字节，float 占 4 字节，double 占 8 字节。
    \item union 只需分配足够容纳最大成员的空间即可：\texttt{max(4, 4, 8) = 8} 字节。
    \item \textbf{对比 struct：} struct 的大小 $\geq$ 各成员大小之和（还需考虑对齐填充）。例如 \texttt{struct\{int i; double k;\}} 通常是 16 字节（int 对齐到 8）。
\end{itemize}

\textbf{常见错误：}
\begin{itemize}
    \item 把 union 当 struct 算，把所有成员大小加起来——这是最经典的错误。
    \item 忽略内存对齐因素。实际上 union 在某些平台上也可能因对齐而稍大于最大成员，但本题的选项和常见计算机条件下答案就是 8。
\end{itemize}

\textbf{讲师提示：} \textit{可以用”酒店房间”类比：struct 是每人一间房，union 是所有客人共用一间最大的房。}

{\color{highlight}答案：B（8）}

% -------------------------------------------------------
\subsection{第16题：结构体指针成员访问}

\textbf{概念概要：} 结构体数组、结构体指针、成员访问运算符（\texttt{.} 与 \texttt{->}）的正确用法。

\medskip
\textbf{讲师提示：} \textit{此题考查细节，要让学生理解”数组名”和”结构体变量”的本质区别。}

\textbf{代码回顾：}
\begin{lstlisting}[language=C]
struct student { int num; int age; };
struct student stu[3] = {{1001,20},{1002,19},{1004,20}};
struct student *p;
p = stu;
\end{lstlisting}

\textbf{逐选项分析：}
\begin{itemize}
    \item A. \texttt{(*p).num}：解引用 p 得到结构体，再用 . 取成员。等价于 \texttt{p->num}。\textbf{正确。}
    \item B. \texttt{(p++)->num}：p 先参与 -> 运算取 num，然后 p 自增指向 stu[1]。\textbf{正确。}
    \item C. \texttt{p = stu.age}：\texttt{stu} 是数组名（类型为 \texttt{struct student[3]}），\underline{数组没有 .age 成员}！只有结构体变量/指针才能用 . 或 -> 取成员。\textbf{错误。}
    \item D. \texttt{p++}：指针自增，p 从 stu[0] 移到 stu[1]。\textbf{正确。}
\end{itemize}

\textbf{常见错误：}
\begin{itemize}
    \item 把数组名和结构体变量混淆——stu 是数组，stu[0] 才是结构体。
    \item 试图对数组名使用 .age——这是最初级的语法错误。
    \item 不清楚 \texttt{p++} 移动的距离：p 是 \texttt{struct student *}，++ 会跳过一整个结构体的大小。
\end{itemize}

{\color{highlight}答案：C}

% -------------------------------------------------------
\subsection{第17题：for 循环条件分析}

\textbf{概念概要：} 赋值运算符 \texttt{=} 与相等判断 \texttt{==} 的区别——C 语言中最经典也最危险的陷阱。

\medskip
\textbf{讲师提示：} \textit{这是高考频考点！几乎每届期末都会有一道题考 = 和 == 的混淆。}

\textbf{代码分析：}
\begin{lstlisting}[language=C]
for (j = 0, k = -1; k = 1; j++, k++)
    printf("*****\n");
\end{lstlisting}

\textbf{关键点：}
\begin{itemize}
    \item for 的第二个表达式 \texttt{k = 1} 是\textbf{赋值}，不是判断。在 C 中 \texttt{=} 运算符的返回值是\underline{被赋的值}。
    \item 赋值表达式 \texttt{k = 1} 的值为 1，在 C 中非 0 即为真（true）。
    \item 因此循环条件\textbf{永远为真}——这是无限循环。
    \item 如果写成 \texttt{k == 1}，则初始 k=-1，条件为假，循环体一次也不执行。
\end{itemize}

\textbf{常见错误：}
\begin{itemize}
    \item 把 \texttt{k = 1} 读成”判断 k 是否等于 1“——这是视觉上的条件反射错误。
    \item 不清楚”赋值表达式有值”的概念——在 C 中，a = 5 这个表达式本身的值就是 5。
    \item 不知道 \texttt{j = 0, k = -1} 是逗号表达式：依次求值，最终值为 -1，但 for 的初始化部分只关心副作用（j=0, k=-1）。
\end{itemize}

\textbf{延伸：} 怎么避免这种错误？如果编译器不说，怎么办？——教学生把常量写在左边：\texttt{1 == k}，一旦写成 \texttt{1 = k} 编译器会报错（因为不能给常量赋值）。但这个技巧在考试中不帮你解题，关键是\textbf{仔细读代码}。

{\color{highlight}答案：B（无限循环）}

% -------------------------------------------------------
\subsection{第18题：存储类别 static}

\textbf{概念概要：} static 关键字对全局变量的作用——限制作用域（内部链接），不影响生命周期。

\textbf{核心讲解要点：}
\begin{itemize}
    \item 题干关键词：”\textbf{只允许本源文件所有函数使用}”——这描述的是 static 修饰的全局变量。
    \item static 修饰全局变量/函数 → \textbf{内部链接}（internal linkage）：仅在当前 .c 文件内可见，其他文件无法 extern 引用。
    \item 对比：不加 static 的全局变量默认是\textbf{外部链接}（external linkage），其他文件可通过 \texttt{extern} 声明后访问。
    \item auto：局部变量默认存储类别，离开作用域即销毁——与”全局”矛盾。
    \item register：建议编译器将变量放入寄存器（仅局部变量）——与”全局”矛盾。
    \item extern：声明外部变量，意在跨文件使用——与”只允许本源文件”矛盾。
\end{itemize}

\textbf{记忆口诀：} static 的两个含义（根据位置不同）：
\begin{enumerate}
    \item 修饰\textbf{局部变量} → 延长生命周期（函数调用间保值）
    \item 修饰\textbf{全局变量/函数} → 限制可见性（仅本文件）
\end{enumerate}
二者本质相通：都是”静态 + 私密”。

\textbf{常见错误：}
\begin{itemize}
    \item 混淆 static 修饰局部变量和全局变量的不同含义。
    \item 认为 static 全局变量”不能修改”——实际上可以在本文件内任意修改，只是外部不可见。
\end{itemize}

{\color{highlight}答案：C}

% -------------------------------------------------------
\subsection{第19题：字符串指针数组——类型判断}

\textbf{概念概要：} 二维字符数组与指针的类型推导。

\textbf{代码分析：}
\begin{lstlisting}[language=C]
char s[][81] = {"Student","Teacher","Father","Mother"}, *ps = s[2];
\end{lstlisting}

\textbf{类型推导 — 从外到内：}
\begin{itemize}
    \item \texttt{s}：由初始化列表推断为 \texttt{char[4][81]}——4 行，每行 81 个字符。
    \item \texttt{s[2]}：取第 3 行（索引从 0 开始），类型为 \texttt{char[81]}。
    \item 在初始化表达式中，\texttt{char[81]} \textbf{退化为} \texttt{char *}（指向首元素的指针）。
    \item 因此 \texttt{*ps} 的声明等价于：\texttt{ps} 是 \texttt{char *} 类型，初始化为指向 "Father" 的首字符。
\end{itemize}

\textbf{常见错误：}
\begin{itemize}
    \item 选 \texttt{char[81]}——没有理解数组退化为指针的规则。
    \item 选 \texttt{char **}——s[2] 是一维数组退化为指针，不是”指针的指针”。\texttt{char **} 对应的是 \texttt{s} 在某些上下文中的退化（s 退化为 \texttt{char (*)[81]}，也不是 char **）。
\end{itemize}

\textbf{讲师提示：} \textit{强调”数组名在表达式中退化为指向首元素的指针”——这是 C 语言最核心的规则之一，贯穿整个指针体系。}

{\color{highlight}答案：B（char *）}

% -------------------------------------------------------
\subsection{第20题：字符串指针输出}

\textbf{概念概要：} 指针解引用的多层含义——\texttt{*ps} vs \texttt{ps} vs \texttt{*s[1]}。

\medskip
\textbf{讲师提示：} \textit{结合第 19 题的上下文，此题需要把 s[1]、ps、*ps 三者的输出格式串讲清楚。}

\textbf{代码（基于第19题）：}
\begin{lstlisting}[language=C]
printf("%c,%s,%c\n", *s[1], ps, *ps);
\end{lstlisting}

\textbf{逐项分析：}
\begin{itemize}
    \item \texttt{*s[1]}：\texttt{s[1]} 是 \texttt{"Teacher"}（类型 char[81]，退化为 char*），\texttt{*} 解引用取首字符，得 \texttt{'T'}。格式符 \texttt{\%c}，输出 \texttt{T}。
    \item \texttt{ps}：指向 \texttt{s[2]} = \texttt{"Father"} 的首字符。格式符 \texttt{\%s} 要求 char*，输出整个字符串 \texttt{Father}。
    \item \texttt{*ps}：解引用 ps，取 \texttt{"Father"} 的首字符，得 \texttt{'F'}。格式符 \texttt{\%c}，输出 \texttt{F}。
\end{itemize}

\noindent 最终输出：\texttt{T, Father, F}

\textbf{常见错误：}
\begin{itemize}
    \item 把 \texttt{*s[1]} 当成 “Teacher“——忘记 \texttt{*} 解引用取的是单个字符。
    \item 把 \texttt{*ps} 当成 “Father“——同样的错误，\%c 输出的是字符而非字符串。
    \item 混淆 \texttt{\%s} 和 \texttt{\%c} 对参数的要求：\%s 需要指针（地址），\%c 需要字符值。
\end{itemize}

{\color{highlight}答案：A（T, Father, F）}
